\documentclass[11pt,twocolumn]{article}

%% ── Packages ─────────────────────────────────────────────────────────────────
\usepackage[T1]{fontenc}
\usepackage[utf8]{inputenc}
\usepackage{lmodern}
\usepackage[margin=1in,columnsep=0.35in]{geometry}
\usepackage{amsmath,amssymb}
\usepackage{graphicx}
\graphicspath{{results/}}
\usepackage{booktabs}
\usepackage{multirow}
\usepackage{array}
\usepackage{xcolor}
\usepackage{hyperref}
\usepackage{cite}
\usepackage{caption}
\usepackage{subcaption}
\usepackage{float}
\usepackage{microtype}
\usepackage{setspace}
\usepackage{url}
\usepackage{siunitx}
\usepackage{threeparttable}
\usepackage{rotating}
\usepackage{pdflscape}
%% tikz is not used — flowchart is SVG-based; do not add tikz.

%% ── Hyperref ─────────────────────────────────────────────────────────────────
\hypersetup{
  colorlinks=true, linkcolor=blue!70!black,
  citecolor=green!50!black, urlcolor=blue!70!black,
}

%% ── Custom commands ──────────────────────────────────────────────────────────
\newcommand{\best}[1]{\textbf{#1}}
\newcommand{\enc}[1]{\texttt{#1}}
\definecolor{bestcol}{HTML}{27AE60}
\newcommand{\bestval}[1]{{\color{bestcol}\textbf{#1}}}

%% ── Title and authors ────────────────────────────────────────────────────────
\title{\textbf{Predicting Surgical Case Duration with Lightweight Language Model Embeddings:
A LangChain--LangGraph Pipeline Combining Text, Structured, and LLM-Extracted Features}}

\author{%
  [Author Names]\\
  \small Department of [Surgery / Biomedical Informatics], [Institution]\\
  \small \texttt{[email@institution.edu]}
}
\date{April 2026}

%% =============================================================================
\begin{document}
\maketitle

%% ── Abstract ─────────────────────────────────────────────────────────────────
%% One continuous paragraph, max 250 words.
%% Structure: Background → Methods → Results → Conclusions.
%% No \cite{}, \ref{}, or undefined abbreviations.
\begin{abstract}
\textbf{Background:} Accurate prediction of surgical case duration is critical for
operating-room scheduling, resource allocation, and reduction of costly schedule
overruns. Existing models rely predominantly on structured tabular features, leaving
rich semantics in free-text procedure names underutilised.

\textbf{Methods:} We developed a five-stage automated pipeline—orchestrated with
LangGraph—that integrates (i) structured electronic health-record features,
(ii) 384-dimensional sentence embeddings of procedure names (HuggingFace
\texttt{all-MiniLM-L6-v2}), and (iii) structured clinical attributes extracted by a
large language model (Groq Llama-3.1-8b-instant) including procedural complexity,
body region, and surgical modality.  Four encoding variants were evaluated across
seven machine-learning models (Ridge, Lasso, ElasticNet, Random Forest, XGBoost,
LightGBM, MLP) with Optuna hyperparameter optimisation and rigorous 5-fold
cross-validation on 180,370 surgical cases.

\textbf{Results:} The best model—LightGBM trained on the combined
structured + HuggingFace embedding + LLM-feature encoding—achieved a mean absolute
error (MAE) of \bestval{26.58~min} ($\pm$0.11), RMSE of 42.26~min, R$^2$ of 0.852,
and SMAPE of 21.69\%.  Adding HuggingFace embeddings to structured features alone
reduced MAE by 26\% (46.6 $\to$ 34.1~min for Random Forest) and by 32\%
(35.0 $\to$ 26.8~min for LightGBM).  LLM-extracted features provided a further
consistent 0.2--0.5~min improvement over embeddings alone.

\textbf{Conclusions:} Lightweight text embeddings and LLM-extracted procedural
features substantially improve surgical case-time prediction beyond structured data
alone.  The fully reproducible open-source pipeline achieves state-of-the-art
performance using only free-tier API services and locally-run models.

\medskip
\noindent\textbf{Keywords:} surgical case duration, operating room scheduling,
natural language processing, sentence embeddings, LightGBM, LangChain, LangGraph,
machine learning.
\end{abstract}

%% ── Abbreviations ─────────────────────────────────────────────────────────────
%% Update whenever a new acronym is introduced anywhere in the paper.
\section*{Abbreviations}
\begin{description}
    \item[ASA]          American Society of Anesthesiologists
    \item[EHR]          Electronic health records
    \item[HF]           HuggingFace
    \item[HPO]          Hyperparameter optimisation
    \item[IQR]          Interquartile range
    \item[LightGBM]     Light Gradient Boosting Machine
    \item[LLM]          Large language model
    \item[MAE]          Mean absolute error
    \item[MLP]          Multilayer perceptron
    \item[NLP]          Natural language processing
    \item[OR]           Operating room
    \item[PCA]          Principal component analysis
    \item[RMSE]         Root mean squared error
    \item[SMAPE]        Symmetric mean absolute percentage error
    \item[SQLite]       Self-contained SQL database engine
    \item[TPE]          Tree-structured Parzen Estimator
    \item[XGBoost]      Extreme Gradient Boosting
\end{description}

%% ── 1. Introduction ──────────────────────────────────────────────────────────
\section{Introduction}\label{sec:introduction}
%% Introduction — migrate content from manuscript.tex \section{Introduction}
%% Step 5: list papers needed as PAPER: Exact Title (one per line, no explanation)
\textcolor{gray}{\textit{[Introduction will be written in Step 5.]}}


%% ── 2. Related Work ──────────────────────────────────────────────────────────
\section{Related Work}\label{sec:related_work}
%% Related Work — migrate content from manuscript.tex \section{Related Work}
%% Step 6: list search phrases as SEARCH: phrase (one per line, no explanation)
\textcolor{gray}{\textit{[Related Work will be written in Step 6.]}}


%% ── 3. Methodology ───────────────────────────────────────────────────────────
\section{Methodology}\label{sec:methodology}
%% Methodology — migrate content from manuscript.tex \section{Methods}
%% All tables written inline in this file — never create separate .tex table files.
%% Reference the flowchart as Figure~\ref{fig:flowchart} (defined in 00_main.tex).
\textcolor{gray}{\textit{[Methodology will be written in Step 3.]}}


%% Flowchart referenced inside 03_methodology as Figure~\ref{fig:flowchart}
\begin{figure*}[tp]
  \centering
  \input{flowchart/06_flowchart}
  \caption{Overall workflow of the proposed five-stage LangGraph pipeline.}
  \label{fig:flowchart}
\end{figure*}

%% ── 4. Results and Discussion ────────────────────────────────────────────────
\section{Results and Discussion}\label{sec:results}
%% Results and Discussion — migrate content from manuscript.tex \section{Results}
%% All tables written inline in this file — never create separate .tex table files.
\textcolor{gray}{\textit{[Results and Discussion will be written in Step 4.]}}


%% ── 5. Conclusion ────────────────────────────────────────────────────────────
%% Write conclusion inline here in Step 7. Four paragraphs:
%% 1. Summary — restate problem + approach in 2–3 sentences. No numbers.
%% 2. Contribution — key quantitative findings; what makes this work novel.
%% 3. Limitations — dataset scope, generalisability, assumptions. No hedging.
%% 4. Future work — 2–3 concrete directions tied to limitations or open questions.
\section{Conclusion}\label{sec:conclusion}
\textcolor{gray}{\textit{[Conclusion will be written in Step 7.]}}

%% ── CRediT authorship contribution statement ─────────────────────────────────
\section*{CRediT authorship contribution statement}
\textbf{Mohammad Noorcheshm}: Conceptualization, Methodology, Data analysis
and curation, Software, Writing: original draft and editing.

%% ── Declaration of competing interest ───────────────────────────────────────
\section*{Declaration of competing interest}
The authors declare that they have no known competing financial interests or
personal relationships that could have appeared to influence the work reported
in this paper.

%% ── Acknowledgement ─────────────────────────────────────────────────────────
\section*{Acknowledgement}
%% Fill in funding sources and acknowledgements before submission.

%% ── Data availability ────────────────────────────────────────────────────────
%% Choose one option and remove the others before submission:
%% Option A — data is open:
%%   The data used in this study are publicly available at [repository/DOI].
%% Option B — data is proprietary:
%%   The data used in this study is proprietary and subject to confidentiality
%%   agreements; therefore, it cannot be shared.
%% Option C — data available on request:
%%   The data that support the findings of this study are available from the
%%   corresponding author upon reasonable request.
\section*{Data availability}

%% ── Bibliography ─────────────────────────────────────────────────────────────
\bibliographystyle{unsrt}
\bibliography{overleaf/references}

\end{document}
